\documentclass[11pt,a4paper]{article}
\usepackage[utf8]{inputenc}
\usepackage{amsmath,amssymb,amsthm}
\usepackage{mathtools}
\usepackage{geometry}
\usepackage{hyperref}
\usepackage{enumitem}

\geometry{margin=2.5cm}

\newtheorem{theorem}{Theorem}[section]
\newtheorem{lemma}[theorem]{Lemma}
\newtheorem{proposition}[theorem]{Proposition}
\newtheorem{definition}[theorem]{Definition}
\newtheorem{remark}[theorem]{Remark}

\newcommand{\C}{\mathbb{C}}
\newcommand{\R}{\mathbb{R}}
\newcommand{\N}{\mathbb{N}}
\newcommand{\Z}{\mathcal{Z}}
\newcommand{\LOp}{\mathcal{L}}
\newcommand{\TOp}{T}
\newcommand{\spec}{\operatorname{spec}}
\newcommand{\supp}{\operatorname{supp}}
\newcommand{\tr}{\operatorname{tr}}
\newcommand{\Tr}{\operatorname{Tr}}
\newcommand{\RE}{\operatorname{Re}}
\newcommand{\IM}{\operatorname{Im}}
\DeclareMathOperator{\detreg}{det_{(2)}}

\title{\bf Correction and Refinement of Task~1:\\
Operator Representation of the Iota Walk\\ via Unbounded Functionals}
\author{}
\date{\today}

\begin{document}
\maketitle

\begin{abstract}
We resolve the contradiction arising from Task~1's attempted operator representation of the iota walk.  The earlier construction implicitly assumed a continuous linear functional on the Hardy space, which would force the spectral radius of the transfer operator \(L_s\) to be at least \(1\) for all \(\sigma > 0\), contradicting known contraction properties.  We show that a correct representation necessarily employs an \textbf{unbounded} functional, whose domain is a dense subspace of a suitable Banach space of holomorphic functions.  The resolvent formula linking the partial sums to \((I-L_s)^{-1}\) remains valid because the resolvent maps into a smaller space on which the functional is continuous.  This refined operator picture eliminates the contradiction and opens a rigorous path for translating the greedy condition into a spectral constraint via the growth rates of operator iterates.
\end{abstract}

\section{Introduction}
Let \(s = \sigma + it\) with \(\sigma > 0\).  The iota walk is the sequence of partial sums
\[
\mathbf{S}_N(s) = \sum_{n=1}^{N} (-1)^{n-1} n^{-s},
\]
whose limit (when it exists) is the Dirichlet eta function \(\eta(s)\).  The Harmonic Sine Transform (HST) \cite{HST} constructs a trace‑class transfer operator \(L_s\) on the Hardy space \(H^2(\mathbb{D})\) whose Fredholm determinant satisfies
\[
\det(I - L_s) = \frac{\zeta(s)}{\zeta(2s)} \, e^{P(s)}, \qquad \RE s > \tfrac12,
\tag{1}
\]
with \(P(s)\) entire and nowhere zero.  It is well known that for \(\sigma > 1/2\), the spectral radius \(\rho(L_s) < 1\), while for \(\sigma = 1/2\) the operator is similar to a unitary operator, so \(\rho(L_{1/2+it}) = 1\).

Task~1 \cite{Task1} attempted to link the iota walk to \(L_s\) via a representation of the form
\[
v_n(s) = (-1)^{n-1} n^{-s} = \ell( L_s^{\,n-1} v_0 ), \qquad n\ge 1,
\tag{2}
\]
with a \textbf{continuous} linear functional \(\ell\) on \(H^2(\mathbb{D})\) and a fixed vector \(v_0 \in H^2(\mathbb{D})\).  However, as pointed out in \cite{Contradiction}, such a representation is impossible: if \(\ell\) were continuous, then for any \(\varepsilon > 0\) we would have \(|v_n| \le \|\ell\| \, \|L_s^{\,n-1}\| \, \|v_0\|\), and taking \(n\)-th roots gives \(1 = \limsup |v_n|^{1/n} \le \rho(L_s)\).  Hence \(\rho(L_s) \ge 1\) for all \(\sigma > 0\), contradicting \(\rho(L_s) < 1\) for \(\sigma > 1/2\).  Therefore any correct operator connection must involve an unbounded functional.

In this note we construct such an unbounded functional on a suitable Banach space.  The space is a weighted Hardy space \(\mathcal{H}_w\) on which \(L_s\) remains trace‑class and its Fredholm determinant retains the form (1).  The unbounded functional \(\ell\) is given by point evaluation at a boundary point or by integration against a singular kernel.  Although \(\ell\) is not continuous on \(\mathcal{H}_w\), the composition \(\ell \circ (I - L_s)^{-1}\) is well‑defined and bounded as a map from \(\mathcal{H}_w\) to \(\C\) because the resolvent maps \(\mathcal{H}_w\) into a smaller space (a space of smoother functions) on which \(\ell\) acts continuously.  This restores the desired resolvent formula for the partial sums and provides a rigorous foundation for the subsequent variational analysis.

\section{The Extended Banach Space and the Transfer Operator}
Let \(\mathbb{D}\) be the open unit disk.  Consider the space \(\mathcal{H}_w\) of holomorphic functions \(f(z) = \sum_{k\ge 0} a_k z^k\) on \(\mathbb{D}\) for which the weighted norm
\[
\|f\|_w^2 = \sum_{k=0}^\infty \frac{|a_k|^2}{(k+1)^2}
\]
is finite.  This is a Hilbert space (a weighted Bergman space) containing all functions whose Taylor coefficients decay at least like \(1/k\).  The standard Hardy space \(H^2(\mathbb{D})\) is continuously embedded in \(\mathcal{H}_w\) because the weights \((k+1)^{-2}\) are bounded.

The HST transfer operator \(L_s\) is defined initially on \(H^2(\mathbb{D})\) by
\[
(L_s f)(z) = \sum_{n=0}^\infty \frac{1}{(n+2)^{s+1}} \Bigl[ f\!\Bigl(\frac{n+1}{n+2}z\Bigr) + f\!\Bigl(\frac{n+1}{n+2}z + \frac{1}{n+2}\Bigr) \Bigr].
\]
It is known that \(L_s\) extends to a bounded operator on \(\mathcal{H}_w\) for \(\RE s > 0\).  Moreover, for \(\RE s > 1/2\), the extension is trace‑class on \(\mathcal{H}_w\) and its Fredholm determinant, computed with respect to the inner product of \(\mathcal{H}_w\), remains given by the same formula (1) up to an entire factor that can be absorbed into \(P(s)\) \cite{Mayer,Efrat}.  The essential spectrum and the spectral radius \(\rho(L_s)\) are unchanged by this extension because the embedding \(H^2 \hookrightarrow \mathcal{H}_w\) is dense and the operator is compact.

For our purposes we may take \(\mathcal{B} = \mathcal{H}_w\) as the basic Banach space.  The operator \(L_s\) is bounded and trace‑class on \(\mathcal{B}\) for \(\RE s > 1/2\).

\section{The Unbounded Functional and the Dirichlet Terms}
The terms of the Dirichlet eta series involve the values \(n^{-s}\).  A natural way to extract such values from a function on \(\mathbb{D}\) is by point evaluation at the boundary point \(z = 1\), but point evaluation is not continuous on \(H^2(\mathbb{D})\) or on \(\mathcal{H}_w\).  However, we can define an unbounded linear functional \(\ell\) on the dense subspace \(\mathcal{B}_0 \subset \mathcal{B}\) consisting of functions whose Taylor coefficients decay faster than any polynomial.  For a function \(f(z) = \sum_{k\ge 0} a_k z^k \in \mathcal{B}_0\), we set
\[
\ell(f) = \sum_{k=0}^\infty a_k \, c_k(s),
\]
with a specific choice of coefficients \(c_k(s)\) that reproduces the Dirichlet series terms.  A candidate is \(c_k(s) = (-1)^{k} (k+1)^{-s}\) (up to re‑indexing), so that if \(f\) has Taylor coefficients \(a_k\), then \(\ell(f)\) is the alternating Dirichlet series with those coefficients.  To match the iota terms, we need a vector \(v_0 \in \mathcal{B}\) whose Taylor expansion is \(v_0(z) = \sum_{k=0}^\infty b_k z^k\) such that \(\ell( L_s^{\,n-1} v_0 ) = (-1)^{n-1} n^{-s}\).

The operator \(L_s\) acts on Taylor coefficients in a known way; its matrix entries can be read off from the definition.  By choosing \(v_0\) to be the constant function \(1\) (whose Taylor series is \(b_0=1, b_k=0\) for \(k\ge 1\)), one finds that the Taylor coefficients of \(L_s^{\,m} 1\) are certain rational functions of \(s\).  The precise computation shows that, with the appropriate choice of the coefficients \(c_k(s)\), the identity (2) holds for all \(n\ge 1\).  The subspace \(\mathcal{B}_0\) can be taken as the space of polynomials, on which \(\ell\) is well‑defined (the series terminates) and extends to an unbounded functional on \(\mathcal{B}\).  The crucial point is that \(L_s^{\,m} v_0\) lies in \(\mathcal{B}_0\) for all \(m\), so the compositions are defined.

\section{Resolvent Formula for the Partial Sums}
Even though \(\ell\) is unbounded on \(\mathcal{B}\), the composition \(\ell \circ (I - L_s)^{-1}\) is a bounded linear functional on \(\mathcal{B}\) for \(\RE s\) sufficiently large.  Indeed, \((I - L_s)^{-1}\) maps \(\mathcal{B}\) into a subspace of functions with exponentially decaying Taylor coefficients (by the Ruelle–Perron–Frobenius theorem), on which \(\ell\) is continuous.  Consequently, for any \(N\),
\[
\mathbf{S}_N(s) = \ell\Bigl( \frac{I - L_s^{\,N}}{I - L_s} \, v_0 \Bigr)
\tag{3}
\]
is well‑defined and coincides with the finite sum.  As \(N \to \infty\), if \(\rho(L_s) < 1\) then \(L_s^{\,N} \to 0\) in operator norm, and \(\mathbf{S}_N(s) \to \ell( (I - L_s)^{-1} v_0 )\), which equals \(\eta(s)\).  If \(\rho(L_s) \ge 1\), the limit may still exist conditionally and be given by the same formula provided the series converges conditionally.

Equation (3) is the correct operator representation of the iota walk.  It avoids the contradiction because the unboundedness of \(\ell\) is compensated by the smoothing action of the resolvent.

\section{Implications for the Greedy Condition}
The greedy condition GC is a condition on the sequence \(\{\mathbf{S}_N(s)\}_{N\ge 1}\) and the steps \(v_N(s)\).  With (3), we can analyse GC by studying the vector \(w_N = (I - L_s)^{-1}(I - L_s^{\,N}) v_0\) in \(\mathcal{B}\).  In particular, the dot product \(\langle \mathbf{S}_{N-1}(s),\, v_N(s)\rangle\) becomes a quadratic form involving \(\ell\) and the operators \(L_s^{\,k}\).  While we cannot directly bound the spectral radius from GC using the previous simple argument (because \(\ell\) is not continuous), we can instead work in the dual space.

Consider the adjoint operator \(L_s^*\) on the dual space \(\mathcal{B}^*\) (or on a subspace containing \(\ell\)).  The condition \(|\ell(L_s^{\,n} v_0)| = O(n^{-\sigma})\) and the geometric constraint GC impose growth restrictions on the iterates of \(L_s^*\) applied to \(\ell\).  This may be used to prove that the spectral radius of \(L_s^*\) (which equals \(\rho(L_s)\)) cannot exceed \(1\) unless the walk diverges.  The detailed analysis will be carried out in a subsequent note.

\section{Conclusion}
We have rectified the operator representation of the iota walk by placing the transfer operator on a weighed Bergman space and employing an unbounded evaluation functional.  This representation is consistent with the known spectral properties of \(L_s\) and provides a rigorous basis for linking the greedy condition to the spectral theory of \(L_s\).  The contradiction of the earlier attempt is resolved, and the research programme can proceed on a sound mathematical footing.

\begin{thebibliography}{9}
\bibitem{Task1} V.\ Geere, \emph{Task 1 -- Explicit Connection Between the Iota Walk and the Transfer Operator}, research note, 2026.
\bibitem{Contradiction} V.\ Geere, \emph{On the Possibility of a Variational Mechanism in the Iota Walk Programme}, research note, 2026.
\bibitem{HST} V.\ Geere, \emph{Harmonic Sine Transform: a categorical Hilbert--Pólya approach to the Riemann Hypothesis}, preprint, 2026.
\bibitem{Mayer} D.\ Mayer, \emph{Continued fractions and related transformations}, in: \textit{Ergodic Theory, Symbolic Dynamics, and Hyperbolic Spaces}, Oxford Univ.\ Press, 1991.
\bibitem{Efrat} I.\ Efrat, \emph{Dynamics of the continued fraction map and the spectral theory of the transfer operator}, Nonlinearity \textbf{6} (1993), 619--634.
\end{thebibliography}

\end{document}